\documentclass[11pt, letterpaper]{article}
\usepackage[utf8]{inputenc}
\usepackage[spanish]{babel}
\usepackage{amsmath, amssymb}
\usepackage{geometry}
\usepackage{tcolorbox}
\usepackage{fancyhdr}
\usepackage{titlesec}
\usepackage{xcolor}
\geometry{top=2.5cm, bottom=2.5cm, left=1.6cm, right=1.6cm}
\setlength{\headheight}{24pt}

% Recuadros de Fórmulas
\tcbset{
    colback=blue!5!white,
    colframe=blue!60!black,
    fonttitle=\bfseries,
    arc=2mm,
    boxrule=0.9pt,
    top=3mm,
    bottom=3mm,
    left=3mm,
    right=3mm
}

% Encabezado
\pagestyle{fancy}
\fancyhf{}
\lhead{\small \textbf{Arquitectura del Computador} \\ Facultad Experimental de Ciencias y Tecnología (FACYT)}
\rhead{\small Formulario -- Capítulo 1 \\ \textbf{Christian Báez C.I. 31.473.355}}
\cfoot{\thepage}
\renewcommand{\headrulewidth}{0.5pt}

% Estilo de las secciones
\titleformat{\section}{\large\bfseries\color{blue!60!black}}{}{0pt}{}[\titlerule]
\titlespacing*{\section}{0pt}{14pt}{8pt}

\begin{document}

\begin{center}
    \Large\textbf{Formulario de Prestaciones y Rendimiento de la CPU} \\
    \vskip 0.2cm
    \small \textit{Referencia: Estructura y diseño de computadores: la interfaz software/hardware (Patterson \& Hennessy)}
\end{center}

\section{1. Definición de Prestaciones y Comparación de Rendimiento}

Las prestaciones (rendimiento) de un computador se definen de manera inversa a su tiempo de respuesta o tiempo de ejecución:

\begin{tcolorbox}[title=Relación Fundamental entre Prestaciones y Tiempo]
$$\text{Prestaciones}_X = \frac{1}{\text{Tiempo de ejecución}_X}$$
\end{tcolorbox}

Si un computador $X$ es más rápido que un computador $Y$, se cumple que $\text{Prestaciones}_X > \text{Prestaciones}_Y$, por lo tanto:
$$\frac{1}{\text{Tiempo de ejecución}_X} > \frac{1}{\text{Tiempo de ejecución}_Y} \implies \text{Tiempo de ejecución}_Y > \text{Tiempo de ejecución}_X$$

\subsection*{Comparación Cuantitativa (``$n$ veces más rápido'')}
Para indicar cuantitativamente cuántas veces es más rápida la máquina $X$ respecto a la máquina $Y$, se calcula el cociente de sus prestaciones:

\begin{tcolorbox}
$$\frac{\text{Prestaciones}_X}{\text{Prestaciones}_Y} = \frac{\text{Tiempo de ejecución}_Y}{\text{Tiempo de ejecución}_X} = n$$
\end{tcolorbox}

\section{2. Factores del Tiempo de Ejecución de la CPU}

El rendimiento de la CPU depende directamente de los ciclos de reloj del programa y del periodo/frecuencia del hardware. Se calcula mediante estas dos expresiones equivalentes:

\begin{tcolorbox}[title=Tiempo de CPU en función del Ciclo de Reloj]
$$\text{Tiempo de ejecución de CPU} = \text{Ciclos de reloj de la CPU} \times \text{Tiempo del ciclo de reloj}$$
\end{tcolorbox}

\begin{tcolorbox}[title=Tiempo de CPU en función de la Frecuencia]
$$\text{Tiempo de ejecución de CPU} = \frac{\text{Ciclos de reloj de la CPU}}{\text{Frecuencia de reloj}}$$
\end{tcolorbox}
\textit{Nota: La frecuencia de reloj es el inverso multiplicativo del tiempo de ciclo ($\text{Frecuencia} = 1 / \text{Tiempo de ciclo}$).}

\section{3. La Ecuación Clásica de las Prestaciones de la CPU}

Al desagregar los ciclos del programa introduciendo el \textbf{Número de instrucciones} y los ciclos promedio por instrucción (\textbf{CPI}), se obtiene la ecuación fundamental de la arquitectura:

\begin{tcolorbox}[title=Ecuación Clásica de Rendimiento de la CPU]
$$\text{Tiempo de ejecución} = \text{Número de instrucciones} \times \text{CPI} \times \text{Tiempo de ciclo}$$
$$\text{Tiempo de ejecución} = \frac{\text{Número de instrucciones} \times \text{CPI}}{\text{Frecuencia de reloj}}$$
\end{tcolorbox}

\section{4. Cálculo de Ciclos Totales y CPI Promedio}

Si un programa ejecuta diferentes mezclas de instrucciones donde cada clase posee un CPI específico, los ciclos totales de la CPU se modelan mediante una sumatoria discreta:

\begin{tcolorbox}[title=Ciclos de Reloj Totales (Múltiples Clases)]
$$\text{Ciclos de reloj de la CPU} = \sum_{i=1}^{n} (\text{CPI}_i \times C_i)$$
\end{tcolorbox}
Donde:
\begin{itemize}
    \item $C_i$ es el número de instrucciones ejecutadas pertenecientes a la \textbf{Clase $i$}.
    \item $\text{CPI}_i$ es el promedio de \textbf{Ciclos por Instrucción} específico de la \textbf{Clase $i$}.
\end{itemize}

\subsection*{CPI Promedio Global del Programa}
El CPI global ponderado del programa se determina dividiendo la sumatoria total de ciclos entre el volumen total de instrucciones:
\begin{tcolorbox}
$$\text{CPI} = \frac{\text{Ciclos de reloj de la CPU}}{\text{Número de instrucciones}} = \frac{\sum (\text{CPI}_i \times C_i)}{\text{Total de instrucciones}}$$
\end{tcolorbox}

\section{5. Descomposición Dimensional (Idea Clave)}

Para comprobar la consistencia formal de las métricas, se utiliza el análisis dimensional de unidades para expresar el tiempo total en segundos por programa:

\begin{tcolorbox}[title=Análisis de Unidades del Tiempo de Ejecución]
$$\text{Tiempo} = \frac{\text{Segundos}}{\text{Programa}} = \frac{\text{Instrucciones}}{\text{Programa}} \times \frac{\text{Ciclos de reloj}}{\text{Instrucción}} \times \frac{\text{Segundos}}{\text{Ciclo de reloj}}$$
\end{tcolorbox}

\end{document}
